\documentclass[11pt,a4paper]{article}

\usepackage[utf8]{inputenc}
\usepackage[T1]{fontenc}
\usepackage{lmodern}
\usepackage{microtype}
\usepackage{amsmath,amssymb,amsthm}
\usepackage{mathtools}
\usepackage{hyperref}
\usepackage{geometry}
\geometry{a4paper,margin=25mm}

\newtheorem{theorem}{Theorem}[section]

\title{\textbf{Identity Bedrock: A Minimal Cryptographic Substrate and Its Systemic Implications}}
\author{Jean-François Bouchard}
\date{March 2026}

\begin{document}

\maketitle

\begin{abstract}
We present Identity Bedrock, a deployed identity system constructed from a minimal cryptographic substrate. Starting from the requirement that no institutional key exists beneath the individual, we derive an architecture grounded in one-way functions (OWFs), which are necessary and sufficient for existentially unforgeable signature schemes under adaptive chosen-message attack.

The system separates identity cryptography from directory consensus via a chain-agnostic, append-only binding layer, enabling portability across heterogeneous substrates while preserving verifiability.

We show that additional system capabilities---including naming, structured representation, execution, connectivity, and economic interaction---emerge as derived properties of this identity substrate rather than independent subsystems.

The contribution is not a new primitive but a binding of existing ones into a system that minimizes trust assumptions while preserving verifiable authentication.
\end{abstract}

\section{Introduction}

Digital identity systems commonly rely on institutional trust anchors such as certificate authorities or identity providers, introducing external control points.

This work explores whether identity can be constructed by minimizing such dependencies, using only necessary cryptographic assumptions, and examines the broader system implications of such a construction.

\subsection{Contributions}

\textbf{Minimal construction.} Identity grounded in OWF-based signatures.

\textbf{Chain-agnostic architecture.} Decoupled identity and directory layers.

\textbf{Deployed system.} Publicly verifiable implementation.

\textbf{Systemic emergence.} Higher-layer capabilities arise from the identity substrate.

\section{Background}

One-way functions (OWFs) are necessary and sufficient for existentially unforgeable signature schemes. This equivalence provides a minimal foundation for authentication.

\section{Architecture}

\subsection{Identity Cryptography}

A user generates a keypair $(sk, pk)$ from a high-entropy seed. Authentication consists of signing a challenge.

\subsection{Directory Substrate}

An append-only directory binds identifiers to public keys, providing global verifiability and finality.

\section{Security}

\begin{theorem}
If the signature scheme is existentially unforgeable under adaptive chosen-message attack and the directory is append-only with finality, then authentication is unforgeable under standard assumptions.
\end{theorem}

\begin{proof}
A successful forgery implies a break of the underlying signature scheme or directory integrity.
\end{proof}

\section{Derived System Properties}

We now examine how additional system capabilities emerge from the identity substrate.

\subsection{Naming and Resolution}

Identity enables decentralized naming via append-only key-value bindings, replacing hierarchical naming authorities with ownership-based resolution.

\subsection{Structured Representation}

Structured objects may be expressed in deterministic formats, enabling consistent parsing, indexing, and decentralized query across nodes.

\subsection{Content Addressing}

Content retrieval is performed via cryptographic hashes, ensuring integrity without trusted delivery channels.

\subsection{Deterministic Execution}

Execution can be defined as a function of input, environment, and command:

\begin{equation}
Execution = Hash(Input, Environment, Command)
\end{equation}

This enables reproducible and verifiable computation.

\subsection{Connectivity}

Overlay networking techniques allow nodes to communicate independently of underlying network constraints, reducing reliance on centralized routing or addressability assumptions.

\subsection{Economic Interaction}

Incentive mechanisms based on participation and resource contribution can operate without centralized validation, leveraging the identity substrate for attribution and verification.

\section{Implementation}

The system is implemented using an append-only name-value substrate anchored to external consensus mechanisms. Identity bindings are publicly verifiable.

\section{Evaluation}

Authentication requires only signature verification and directory lookup, resulting in low computational overhead and minimal latency.

The system has been deployed in practice, demonstrating operational stability over extended periods.

\section{Discussion}

The proposed architecture does not introduce multiple independent subsystems. Instead, it demonstrates that naming, execution, connectivity, and economic behaviors can be understood as consequences of a single identity substrate.

This reduces conceptual complexity while preserving compatibility with existing protocols.

\section{Conclusion}

Identity Bedrock demonstrates that a minimal cryptographic substrate, combined with an append-only directory, is sufficient not only for authentication but also for supporting broader system behaviors.

By deriving system functionality from a unified primitive, the approach reduces reliance on centralized infrastructure while maintaining verifiability and portability.

\end{document}